%! AP Statistics — Unit 2, Lesson 2.6 — Warm-Up — Answer Key
\documentclass[10pt]{article}
\usepackage{apstats-article}
\usepackage{apstats-key}

\begin{document}

\pageheader{Unit 2, Lesson 2.6}{Warm-Up --- Answer Key}
\namedateperiod

\vspace{0.3in}

% ── PART 1 ────────────────────────────────────────────────────────────────────
{\large\bfseries\color{navy} Part 1 \quad Slope-Intercept Form (Algebra Review)}

\vspace{0.12in}
\noindent Recall from algebra: in $y = mx + b$, the slope $m$ gives the change in $y$ for
each one-unit increase in $x$, and $b$ is the value of $y$ when $x = 0$.

\vspace{0.1in}

\begin{enumerate}[leftmargin=1.6em, itemsep=0.22in]

  \item A food truck owner models weekly revenue (dollars) from cups of lemonade sold ($x$) with
        the equation $y = 30 + 2.50x$.

        \begin{enumerate}[label=\textbf{(\alph*)}, leftmargin=1.6em, itemsep=8pt]
          \item What is the slope? \ans{2.50} \quad What is the $y$-intercept? \ans{30}

          \item What does the slope mean in this context?\\[8pt]
                \ansline{For each additional cup of lemonade sold, weekly revenue increases by \$2.50.}

          \item What does the $y$-intercept mean? Is it meaningful in this context?\\[8pt]
                \ansline{When 0 cups are sold, predicted weekly revenue is \$30. This could represent
                a base permit cost; selling zero cups is possible, so the intercept is meaningful here.}
        \end{enumerate}

  \item Recall from Lesson 2.5: the correlation coefficient $r$ describes the
        \ans{direction} and \ans{strength} of a \textit{linear} association.

        \vspace{8pt}
        A scatterplot of hours of sleep ($x$) vs.\ reaction time in milliseconds ($y$) has
        $r = -0.88$. Describe in one sentence what this tells you about the relationship.\\[8pt]
        \ansline{There is a strong negative linear association between hours of sleep and reaction
        time: people who sleep more tend to have faster (lower) reaction times, and the points
        follow the linear pattern closely.}

\end{enumerate}

\vspace{0.28in}

% ── PART 2 ────────────────────────────────────────────────────────────────────
{\large\bfseries\color{navy} Part 2 \quad Setting Up for Regression}

\vspace{0.12in}

\begin{tcolorbox}[colback=greenbg, colframe=greenacc, boxrule=0.8pt, arc=2mm,
    left=4mm, right=4mm, top=3mm, bottom=3mm]
    \small
    \textbf{Think about it:} The correlation $r$ tells you the direction and strength of a
    linear relationship, but it cannot answer:
    \begin{itemize}[leftmargin=1.3em, itemsep=3pt]
      \item \emph{How much} does the response variable change per one-unit increase in the
            explanatory variable?
      \item \emph{What value} would you predict for the response if $x = 7$?
    \end{itemize}
    Today you'll learn the tool that answers both: the \textbf{regression line}.
\end{tcolorbox}

\vspace{0.14in}

\noindent A researcher studying used car sales finds $r = -0.93$ for the relationship between a
car's age in years ($x$) and its resale price in \$1000s ($y$). What does $r$ tell you about the
relationship? What does $r$ \emph{not} tell you?\\[8pt]
\ansline{$r = -0.93$ tells us there is a strong negative linear association: older cars tend to
have lower resale prices, and the points follow a linear pattern closely. But $r$ does not say
how much the price drops per year, nor does it give a specific predicted price for a car of a
particular age.}

\begin{teachernote}
Circulate during Part 2 and listen for students who conflate $r$ with the slope $b$. The whole
point of the warm-up is to surface this gap before introducing $\hat{y} = a + bx$.
\end{teachernote}

\end{document}
